\section{Discussion}
\label{sec:discussion}
\subsection{Prediction across care settings}
In six Israeli general hospitals, INTERVenE-Enc distinguished patients who subsequently developed complications more accurately than ss-STraTS. At day four, the difference in weighted AUPRC was 0.053 (95\% CI 0.036--0.070), and the advantage persisted at each daily cutoff from day two to day six. The models differed in size, so this comparison establishes their relative performance in the configurations studied. The advantage over ss-STraTS extends the earlier finding in MIMIC-IV intensive care records to general inpatient care in another health system \citep{oded2026intervene}. Differences in timing accuracy were more modest. They were clearest for the onset of hyperglycemia, while the most accurate model varied across other complications.

The two settings offered different views of a patient's course. During the first four days, the median number of heart-rate readings per day was two in Kineret and 12.5 in MIMIC-IV. For glucose, it was four in Kineret and 2.75 in MIMIC-IV. A Kineret record could therefore document glycemic changes relatively frequently while leaving hours between observations of vital signs. The encoder performed well under both recording patterns when fitted to local data. Case mix and outcome definitions also differed between the studies, limiting what can be attributed to recording frequency alone.

Within Kineret, a model fitted at five hospitals retained discrimination at the sixth. KB exceeded the local prevalence reference at every site, although weighted AUPRC ranged from 0.431 to 0.539. Differences in case mix, event frequency, and local recording practices matter when interpreting that range. At one hospital, rare events made the aggregate interval unreliable. Because the hospital holdouts included only KB, the evidence for its advantage over other models comes from the pooled comparison.

\subsection{What the clinical knowledge base contributed}
In the ICU study, statistically defined intervals preserved much of the encoder's performance, while the clinical vocabulary made it possible to identify which states and trends contributed to a prediction \citep{oded2026intervene}. In Kineret, the choice of intervals also mattered for discrimination. Clinical intervals improved weighted AUPRC by 0.060 (95\% CI 0.046--0.076) over SAX in the same encoder architecture. Much of the gain came from glycemic outcomes, for which the knowledge base contains detailed definitions, and the advantage persisted when each outcome received equal weight.

SAX indicates which statistical band a variable occupied over a given period. A clinical interval can express that glucose remained elevated, that it was falling, or that a sequence of observations met the criteria for a clinical event. The knowledge base was developed for inpatient diabetes care and mapped to the Kineret hospitals' measurement and treatment codes \citep{oded2027kineret-jbi}. This close fit to the clinical question may help explain its advantage. Thresholds, timing rules, vocabulary, and access to derived events all changed together, however. The experiment measures their combined value and leaves open how much each contributed or how well the same approach would work in another specialty.

\subsection{What guideline assessment added}
The clinical study associated adherence to particular recommendations with subsequent outcomes \citep{oded2027kineret-jbi}. In the present study, each guideline score summarized measurements and treatments already represented in the clinical history. Making those relationships explicit improved discrimination for subsequent hyperglycemia and reduced it for mortality. The overall difference was small and uncertain, and macro AUPRC was essentially the same with and without QA. A recommendation can therefore be clinically relevant even when an assessment of adherence adds little to what the model can infer from the record.

The hypoglycemia findings illustrate what the clinical vocabulary makes visible. A raw-observation model can assign a contribution to a glucose reading at a particular time. The interval model can assign one to a fall in glucose or to the timing of insulin after an elevated reading. A falling glucose trend contributed toward higher predicted hypoglycemia risk, and an insulin-after-hyperglycemia pattern scored false contributed toward lower risk. These directions are consistent with the clinical study, in which greater adherence to that recommendation was associated with fewer episodes of severe hyperglycemia and more episodes of hypoglycemia \citep{oded2027kineret-jbi}. They also fit the clinical rationale for coordinating glucose monitoring, insulin, and nutrition in hospital \citep{ada2026hospital}. The pattern score alone cannot resolve the dose, actual intake, or reason for a treatment decision.

For mortality, the direction of an individual contribution and the overall change in discrimination are especially informative together. Routine glucose monitoring contributed toward lower predicted risk, consistent with the association with 30-day mortality in the clinical study. Yet the model with all QA inputs was less accurate at identifying admissions with a death-coded terminal marker in the forecast window. An attribution describes how the fitted model used a pattern. It cannot tell us what would happen if clinicians monitored glucose more often.

\subsection{Interpreting a contribution to predicted risk}
A name such as \emph{decreasing glucose} comes from the clinical rule that defined the interval. The embedding analysis showed that the encoder retained relationships within that vocabulary, with the two boundaries of the same concept closer together than concepts from different clinical families. In a separate test, modest noise or removal of measurements left most covered time with the same numeric state or trend label. The stability of QA patterns and downstream risk estimates remains untested.

Identifying a clinically recognizable input is only part of explaining a prediction. The transcoder estimates sensitivity at occurrences of that input through a sparse approximation to each layer's output. Its reconstruction explained part of the variance on held-out data ($R^2=0.396$--$0.514$). Attaching the transcoders preserved the original predictions across 300 admissions, but the unexplained variation limits how fully the attributions can account for an individual prediction.

\subsection{Limitations and clinical implications}
\label{sec:discussion-limits}
The day-four results apply to adults with diabetes or early hyperglycemia whose records extended beyond four days. Shorter admissions were outside that evaluation. A recorded onset also depends on when clinicians measured the patient or documented an event. We applied the date restriction to an evaluation of previously fitted models, and the intervals quantify sampling variation across patients for those models. Rare events limit the precision of some hospital and subgroup estimates.

Mortality timing has a further limitation. For a death recorded after discharge, the model predicts the time of the death-coded visit end, so its onset error cannot be interpreted as error in predicting the actual time of death.

Guideline scores reflect what can be reconstructed from the record. For example, meal-related rules used assumed meal times, and the summaries assigned each assessment to its recorded start. Before using these scores during care, it would be necessary to establish when all the observations needed to calculate them would have been available. Illness severity and the circumstances prompting treatment can also influence both recorded adherence and subsequent risk. The agreement between selected attributions and the clinical study must be read in that context. Both analyses draw on the same records.

For a clinician deciding whether to observe a patient more closely, the timing of a complication matters alongside its likelihood. Comparisons with training-median predictions are still needed to interpret the absolute errors for both onset and length of stay. For predicted probabilities, the clinical interval model showed closer agreement with observed outcomes than ss-STraTS in the pooled test set (mean expected calibration error 0.0122 and 0.1268, respectively). Choosing a threshold for action would require checking calibration at the intended hospital and weighing missed events against false alerts.
